%! TeX program = xelatex
\documentclass[10pt]{article}

\usepackage[letterpaper,landscape,margin=1.2cm]{geometry}
\usepackage{fontspec}
\usepackage[spanish]{babel}
\usepackage{booktabs}
\usepackage{array}
\usepackage{tabularx}
\usepackage{xcolor}
\usepackage{colortbl}

\setmainfont{Liberation Sans}
\pagestyle{empty}
\setlength{\parindent}{0pt}
\setlength{\arrayrulewidth}{0.4pt}
\arrayrulecolor{gray!65}
\renewcommand{\arraystretch}{1.28}
\newcolumntype{L}[1]{>{\raggedright\arraybackslash}p{#1}}
\newcolumntype{Y}{>{\raggedright\arraybackslash}X}

\begin{document}

{\Large\bfseries Programa de 10 semanas para cierre del proyecto}\par
\vspace{0.3cm}
{\small Se trabaja con fotografías, geolocalizaciones, mapas y relatos existentes.}\par
\vspace{0.5cm}

\begin{tabularx}{\linewidth}{|L{0.08\linewidth}|Y|Y|Y|}
\hline
\textbf{Semana} & \textbf{Objetivo} & \textbf{Actividades principales} & \textbf{Entregable} \\
\hline
1 & Ordenar información y datos faltantes & Definir C-01 a C-04; reunir fotos, relatos, ubicación, uso, materiales, estado y permisos de publicación. & Tabla completa de la información y carpeta organizada de evidencias. \\
\hline
2 & Cerrar evidencia existente & Revisar fotografías actuales e históricas ya disponibles; confirmar nombres locales como tepetate; definir qué se puede publicar. & Archivo de evidencias cerrado y criterios de privacidad definidos. \\
\hline
3 & Preparar cartografía & Exportar mapas desde QGIS: contexto regional, comunidades, puntos generalizados, recorridos y relación con materiales. & Paquete cartográfico preliminar para insertar en ICR. \\
\hline
4 & Redactar diagnóstico C-01 y C-02 & Llenar fichas narrativas y descriptivas; integrar fotos, materiales, estado, cambios y relato asociado. & Fichas completas C-01 y C-02. \\
\hline
5 & Redactar diagnóstico C-03 y C-04 & Completar fichas; tratar órganos como cerramiento vivo o semivivo; documentar bajareque con entramado, relleno y recubrimiento. & Fichas completas C-03 y C-04. \\
\hline
6 & Comparación, fichas, mapas y criterios. & Elaborar cuadro comparativo; sintetizar materiales, usos, estados y cambios; afinar fichas, mapas y criterios. & Revision final de cápitulos correspondientes. \\
\hline
7 & Ajustar introducción, objetivos y metodología & Revisar coherencia entre preguntas, objetivos, metodología, diagnóstico y producto operativo. & Introducción y metodología alineadas con resultados reales. \\
\hline
8 & Redactar conclusiones y revisar marco & Redactar conclusiones por pregunta de investigación; revisar que el marco teórico no prometa más de lo que el diagnóstico entrega. & Conclusiones completas y marco ajustado. \\
\hline
9 & Edición integral & Revisar estilo, repeticiones, citas, tablas, pies de figura, numeración, referencias, anexos y privacidad. & Borrador final completo. \\
\hline
10 & Cierre técnico y entrega & Compilar, corregir warnings relevantes, revisar PDF, generar versión final y respaldos. & PDF final listo para entrega. \\
\hline
\end{tabularx}

\end{document}
